\documentclass{situationist-preprint}

\title{The Limits of Bad Faith}
\msid{PHE-10-b}
\msfamily{PHE-10}
\msdiscipline{Phenomenology}
\mstype{Consequence}
\mskeywords{bad faith; oppression; facticity; transcendence; Sartre; Beauvoir}

\begin{document}
\maketitle

\begin{abstract}
Sartre's notion of bad faith is often thought to license a harsh verdict on the oppressed: if freedom is absolute, whoever says that she cannot escape her condition is lying to herself. We derive from PHE-10-a that the verdict is mistaken. Bad faith is a misrepresentation of one of two distinct relations, between facts and their meaning and between facts and outcomes. An oppressed person who reports the limits of what she can bring about misrepresents neither, and one to whom oppression has hidden certain possibilities is not lying to herself about them. The charge of bad faith in cases of oppression applies principally to those who present the situation of the oppressed as fixed by their nature.
\end{abstract}

\section{Bad faith as misrepresentation}

Sartre describes bad faith as a lie to oneself that plays on the double character of human reality: it presents transcendence as if it were facticity, or facticity as if it were transcendence \citep[Part~I, ch.~2]{Sartre1943}. PHE-10-a gives the two terms a precise sense. Ontological freedom is the absence of any fact that fixes, by itself, its meaning for the agent; practical freedom is the range of outcomes she can bring about (Definition~2.3 of PHE-10-a). Ontological freedom is absolute (Lemma~3.1), practical freedom is limited and graded (Lemma~3.2), and the two are consistent because they concern different relations (Theorem~3.3). Bad faith can therefore be stated as follows.

\begin{definition}\label{d:bf}
An agent is in \emph{bad faith} about a fact $F$ if, while in a position to recognize the truth, she represents $F$ as fixing its own meaning for her (denial of ontological freedom) or represents an outcome that $F$ excludes as one she can bring about (denial of facticity).
\end{definition}

The clause \enquote{in a position to recognize the truth} expresses Sartre's requirement that bad faith be a lie, not an error: the deceiver and the deceived are one, and what is hidden is in some sense known.

\section{Three consequences}

\begin{proposition}\label{p:report}
A true report of the limits of one's practical freedom is not bad faith, however narrow those limits are.
\end{proposition}

\begin{proof}
By Lemma~3.2 of PHE-10-a, practical freedom is limited by facticity. To report that an outcome is excluded, when it is, is to represent the relation between facts and outcomes correctly, and it represents no fact as fixing its meaning. Neither clause of Definition~\ref{d:bf} is met.
\end{proof}

The worker who says she cannot leave her job, because she would not eat, is not in bad faith.

\begin{proposition}\label{p:hidden}
An agent to whom a possibility has been hidden by the conduct of others is not in bad faith in failing to envisage it.
\end{proposition}

\begin{proof}
Bad faith requires that the agent be in a position to recognize what she conceals from herself. As the reply to Objection~4 of PHE-10-a states, oppression can narrow not only what one can do but what possibilities appear to one; this is a limitation of practical freedom, including the freedom to form projects. Where the possibility of revolt does not appear to her, there is nothing she knows and hides from herself.
\end{proof}

This is Beauvoir's verdict on those kept in ignorance and dependence, whose situation is closed to them by others: they are not to be reproached as the person who flees a freedom she has in view is reproached \citep{Beauvoir1947}.

\begin{proposition}\label{p:oppressor}
Whoever represents the situation of the oppressed as fixed by their nature, while in a position to recognize that it is not, is in bad faith.
\end{proposition}

\begin{proof}
To represent a social position as fixed by a nature is to represent facts (of sex, race or birth) as fixing what they mean for those who bear them. By Lemma~3.1 of PHE-10-a, no fact fixes its own meaning for any agent, and the first clause of Definition~\ref{d:bf} applies to whoever maintains the representation against what she is in a position to recognize.
\end{proof}

Sartre's portrait of the anti-Semite, who chooses to regard the Jew as a fixed essence so that he himself may be a thing among things, is the paradigm \citep{Sartre1946}; Beauvoir's analysis of the myth of the eternal feminine is another \citep{Beauvoir1949}. The same clause applies to the oppressed only where they adopt that representation of themselves while able to see through it, and this is the narrow residue of the charge.

\begin{objection}
If ontological freedom is absolute, the oppressed remain responsible for the meaning they give their situation, and the harsh verdict returns at that level.
\end{objection}

\begin{reply}
Ontological freedom concerns the relation between facts and meaning, not the range of meanings available; which projects she can form is a matter of practical freedom, which others' conduct limits (Proposition~\ref{p:hidden}). Responsibility for taking up a meaning presupposes that it was available; Lemma~3.1 of PHE-10-a does not say that it always is.
\end{reply}

\section{Conclusion}

In cases of oppression, bad faith does not reach truthful reports of constraint or failures to envisage what has been concealed; it reaches, above all, those who present a situation they maintain as a nature.

\begin{thebibliography}{9}\small
\bibitem[Beauvoir(1947)]{Beauvoir1947} Beauvoir, S. de (1947). \emph{Pour une morale de l'ambigu\"it\'e}. Gallimard. Trans. B. Frechtman, \emph{The Ethics of Ambiguity}, Philosophical Library, 1948.
\bibitem[Beauvoir(1949)]{Beauvoir1949} Beauvoir, S. de (1949). \emph{Le Deuxi\`eme Sexe}. Gallimard.
\bibitem[Sartre(1943)]{Sartre1943} Sartre, J.-P. (1943). \emph{L'\^Etre et le n\'eant}. Gallimard.
\bibitem[Sartre(1946)]{Sartre1946} Sartre, J.-P. (1946). \emph{R\'eflexions sur la question juive}. Paul Morihien.
\end{thebibliography}

\end{document}
